\documentclass[10pt,twocolumn]{article}
\usepackage[T1]{fontenc}
\usepackage[utf8]{inputenc}
\usepackage{lmodern}
\usepackage[margin=1.8cm,top=2.4cm,bottom=2cm,columnsep=0.6cm]{geometry}
\usepackage{fancyhdr}
\usepackage{titlesec}
\usepackage{booktabs}
\usepackage{amsmath,amssymb,amsfonts}
\usepackage{microtype}
\usepackage{xcolor}
\usepackage{mdframed}
\usepackage{parskip}
\usepackage{enumitem}
\usepackage{hyperref}

\definecolor{rulered}{RGB}{180,30,30}
\definecolor{boxgray}{RGB}{245,245,245}
\definecolor{keywordgray}{RGB}{100,100,100}

\pagestyle{fancy}
\fancyhf{}
\fancyhead[L]{\textsc{\small Claude Mag}}
\fancyhead[C]{\small\textit{Machine Learning Research Digest}}
\fancyhead[R]{\small 2026-09-09}
\fancyfoot[C]{\small\thepage}
\renewcommand{\headrulewidth}{0.8pt}

\titleformat{\section}{\normalsize\bfseries\color{rulered}}{}{0pt}{}%
  [\vspace{-4pt}\color{rulered}\rule{\columnwidth}{0.4pt}]
\titlespacing{\section}{0pt}{8pt}{4pt}

\hypersetup{colorlinks=true,urlcolor=rulered,linkcolor=black}

\begin{document}
\setlength{\parindent}{0pt}
\setlength{\parskip}{4pt}

{\small\color{keywordgray}\textbf{Keywords:}
  vision-language models \textbullet{}
  causal intervention \textbullet{}
  cross-modal attention \textbullet{}
  video understanding}\par\vspace{4pt}

{\LARGE\bfseries\raggedright
  From Vision to Language: Investigating Causal
  Information Flow in Multimodal
  Decision-Making}\par\vspace{4pt}

{\large\itshape\raggedright
  Causal Patching of Video-Text Attention Reveals
  That VLMs Ground Vision at Answer Tokens,
  with Nouns as Anchors and Temporal Reasoning Fragile}\par\vspace{6pt}

{\small\textbf{By} Davide Testa, Hugh Mee Wong, Alessandro Lenci,
  Bernardo Magnini, Albert Gatt
  (FBK; Univ.\ Malta; Univ.\ Pisa; et al.)}\par\vspace{2pt}

{\small\color{keywordgray}arXiv:2609.05149 \textbullet{} Preprint, September 2026}
\par\vspace{2pt}\noindent{\color{rulered}\rule{\columnwidth}{0.6pt}}\vspace{6pt}

Vision-language models (VLMs) are widely deployed for video
understanding tasks --- question answering, visual reasoning, action
recognition. Yet the internal mechanism by which visual information is
integrated into language-based decisions remains largely opaque.
This paper applies layer-wise causal intervention on video-text attention
pathways to answer a basic question: at which layer, and at which token
positions, does visual information causally determine the model's final
answer? The answer is consistent across five VLMs spanning two orders of
magnitude in parameters: visual information is integrated primarily when
the model processes the candidate \emph{answer} tokens, not the question.
Nouns act as semantic anchors for cross-modal grounding; verbs carry
temporal relational load; and temporal reasoning exhibits a distinct
fragility suggesting models rely on linguistic cues rather than genuine
frame-by-frame sequential understanding.

\begin{mdframed}[backgroundcolor=boxgray,linecolor=rulered,linewidth=1pt,
    innerleftmargin=6pt,innerrightmargin=6pt,
    innertopmargin=6pt,innerbottommargin=6pt]
\textbf{\small AT A GLANCE}\par\vspace{4pt}
\begin{description}[leftmargin=0pt,labelsep=4pt,itemsep=2pt,parsep=0pt,
    font=\small\bfseries]
  \item[Problem:] Where and when do VLMs integrate visual signals when
    answering video questions?
  \item[Why it matters:] Understanding integration sites guides
    architecture choices and reveals failure modes.
  \item[Method:] Layer-wise causal patching of video-text attention
    activations; tested on spatial, causal, and temporal QA.
  \item[Result:] Visual integration peaks at answer tokens; nouns
    anchor cross-modal binding; temporal reasoning fragility TFI=0.34.
  \item[Evidence:] 5 VLMs (7B--72B); consistent across spatial, causal,
    and temporal tasks.
\end{description}
\end{mdframed}

\section{The Question}

VLMs process videos as sequences of frame tokens and answer multiple-choice
questions about their content. The prevailing assumption is that visual
grounding is distributed throughout the forward pass. But mechanistic
interpretability work on image-text models has suggested that
cross-modal integration is often concentrated at specific layers and
token positions. For video, with its additional temporal structure,
the picture is even less clear.

This paper runs the first systematic causal intervention study on the
cross-modal attention pathways of video VLMs, targeting three reasoning
types: spatial (object location, spatial relations), causal (action
outcomes, event causation), and temporal (sequence, duration,
before/after).

\section{Causal Patching Methodology}

The central method is \textbf{causal activation patching} applied to
cross-modal attention activations.

For a clean run on a video-question pair, all attention activations
$\mathbf{A}^{(l)}_t$ at layer $l$, token position $t$ are recorded.
A corrupted run processes a corrupted video (shuffled or masked frames)
to obtain degraded activations.

\textbf{Causal patching} restores the clean activation at a single
$(l^*, t^*)$ while all other activations remain corrupted:
\[
  \mathbf{A}^{(l)}_t =
  \begin{cases}
    \mathbf{A}^{(l)}_t[\text{clean}]     & l = l^*,\; t = t^* \\
    \mathbf{A}^{(l)}_t[\text{corrupted}] & \text{otherwise}
  \end{cases}
\]
The \textbf{causal effect} of visual information at $(l^*, t^*)$ is:
\[
  \text{CE}(l^*, t^*) =
  P(\text{correct}\,|\,\text{patched at } l^*, t^*)
  - P(\text{correct}\,|\,\text{fully corrupted})
\]
High CE means visual information at that position is causally necessary.

Normalising by the maximum recoverable effect:
\[
  \widehat{\text{CE}}(l^*, t^*) =
  \frac{\text{CE}(l^*, t^*)}{P(\text{correct}|\text{clean}) -
                              P(\text{correct}|\text{corrupted})}
\]

\section{Where Visual Integration Happens}

The results are striking in their consistency. Across all five models
and all three task types, the highest causal effects appear at the
positions of \emph{candidate answer tokens} --- not the question,
not the reasoning prefix. The model reads the video primarily when it
evaluates each answer option against the video content.

\begin{center}
\begin{tabular}{lccc}
\toprule
 & Spatial & Causal & Temporal \\
\midrule
Peak CE at answer tokens   & 0.71 & 0.68 & 0.53 \\
Peak CE at question tokens & 0.18 & 0.22 & 0.19 \\
\bottomrule
\end{tabular}
\end{center}

Visual integration peaks at layers 60--75\% of model depth, in the
middle-to-late transformer layers.

\section{Nouns Anchor, Verbs Carry Temporal Load}

Aggregating causal effects by token part-of-speech reveals the role
of different word classes:
\[
  \widehat{\text{CE}}_\text{POS}(\text{pos}) =
  \frac{1}{|T_\text{pos}|} \sum_{t \in T_\text{pos}} \sum_l
  \widehat{\text{CE}}(l, t)
\]

\begin{center}
\begin{tabular}{lccc}
\toprule
POS & Spatial & Causal & Temporal \\
\midrule
Noun share  & 58\% & 51\% & 44\% \\
Verb share  & 12\% & 24\% & 41\% \\
Other share & 30\% & 25\% & 15\% \\
\bottomrule
\end{tabular}
\end{center}

Nouns dominate as binding sites for object identity and spatial
information. Verbs grow in importance for causal and temporal tasks,
where relational and sequential information must be extracted from
the video.

\section{Temporal Fragility}

The \textbf{Temporal Fragility Index} (TFI) measures how much harder
it is to recover visual information for temporal questions:
\[
  \text{TFI} =
  \frac{\widehat{\text{CE}}_\text{spatial}
      - \widehat{\text{CE}}_\text{temporal}}
       {\widehat{\text{CE}}_\text{spatial}}
\]

TFI $= 0.34$ averaged across models. Restoring a clean visual signal
for a temporal question recovers 34\% less accuracy than for a spatial
question, despite identical video content.

Combined with the high verb CE share on temporal tasks, this pattern
suggests VLMs use linguistic temporal markers (verb tense, temporal
adverbs, ``before/after'' phrases) rather than reconstructing sequential
information across frames. Temporal expressions in answer candidates
trigger reliance on linguistic temporal priors.

\section{Experimental Details}

\textbf{Models:} LLaVA-Video-7B, LLaVA-Video-72B, Qwen2.5-VL-7B,
Qwen2.5-VL-72B, InternVL3-8B.

\textbf{Video benchmarks:} Spatial, causal, and temporal questions
derived from standardised evaluation datasets; 4 answer candidates per
question; balanced correct/incorrect.

\textbf{Corruption:} Frame shuffling and frame masking produce similar
CE patterns, confirming that ordering information (not just content
presence) drives temporal fragility.

\textbf{Model scale:} Larger models achieve higher absolute accuracy but
the same qualitative pattern, suggesting these are architectural
properties rather than capacity limitations.

\section*{Reference}

Davide Testa, Hugh Mee Wong, Alessandro Lenci, Bernardo Magnini, and
Albert Gatt. \textbf{From Vision to Language: Investigating Causal
Information Flow in Multimodal Decision-Making.}
arXiv:2609.05149, September 2026.
\url{https://arxiv.org/abs/2609.05149}

\end{document}
